% =====================================================================
%  CAPÍTULO 3 · METODOLOGÍA, PLANIFICACIÓN Y VIABILIDAD
%  Función: explicar CÓMO se ha trabajado y justificarlo.
%  Reglas clave:
%   - Justifica la metodología (por qué esta) y que su nombre corresponda al
%     proceso real («iterativa» solo si de verdad hubo iteraciones).
%   - Un ÚNICO plan: mismo número de fases y nombres en texto y Gantt.
%   - El plan llega hasta el depósito y la defensa, y se redacta en PASADO
%     en la memoria final («se desarrolló», no «se desarrollará»).
%   - Explica las desviaciones: lo planificado que no se hizo se declara.
%   - El stack se presenta UNA sola vez (aquí), con versión exacta y cita.
%   - El Portafoli evalúa la viabilidad económica y legal y la sostenibilidad.
% =====================================================================
\chapter{Metodología, planificación y viabilidad}
\label{cap:metodologia}

Este capítulo describe la metodología seguida, la planificación temporal real
del proyecto, las herramientas utilizadas y el análisis de su viabilidad.

\section{Metodología de desarrollo}
\label{sec:metodologia}

Se ha seguido un desarrollo iterativo e incremental en tres incrementos, cada
uno con análisis, implementación y pruebas propios. Se eligió este enfoque
porque los requisitos de tiempo real dependían de la validación del primer
incremento, y una planificación en cascada habría retrasado ese riesgo hasta el
final del proyecto~\cite{martin2022arquitectura}. Al final de cada incremento se
revisó el avance con la dirección del trabajo.

\section{Planificación temporal}
\label{sec:planificacion}

El proyecto se organizó en tres fases, que se muestran en la
Figura~\ref{fig:cronograma}. La fase de análisis (F1) se desarrolló entre
octubre y diciembre; la de diseño e implementación (F2), entre diciembre y
mayo, y la de validación y cierre (F3), entre mayo y julio, hasta el depósito y
la defensa. Los hitos en rojo corresponden a las entregas que evalúa el
Portafoli de la \ac{EPS}.

\begin{figure}[htbp]
  \centering
  % =====================================================================
%  CRONOGRAMA (diagrama de Gantt) · fuente ÚNICA de la planificación
%  Se usa en la memoria (Cap. 3) y en los informes inicial y de seguimiento:
%  así el plan es el mismo en todos los documentos.
%  - Una columna = una quincena (granularidad uniforme; no la mezcles).
%  - El plan llega SIEMPRE hasta el depósito y la defensa reales.
%  - progress=NN (opcional) marca el avance de una tarea en el seguimiento.
%  - Ajusta los meses y las fechas al calendario oficial del curso.
% =====================================================================
\begin{ganttchart}[
    x unit=0.45cm, y unit chart=0.52cm, y unit title=0.55cm,
    vgrid={*1{draw=black!10}}, hgrid,
    title label font=\scriptsize\bfseries, title height=1,
    bar label font=\footnotesize, group label font=\footnotesize\bfseries,
    milestone label font=\footnotesize\itshape,
    bar/.append style={fill=RoyalBlue!55, draw=RoyalBlue!80!black},
    group/.append style={fill=black!70},
    milestone/.append style={fill=BrickRed, draw=BrickRed},
    progress label text={}, bar incomplete/.append style={fill=RoyalBlue!15},
    bar height=0.55, group height=0.3, milestone height=0.55,
  ]{1}{20}
  \gantttitle{Oct}{2} \gantttitle{Nov}{2} \gantttitle{Dic}{2} \gantttitle{Ene}{2}
  \gantttitle{Feb}{2} \gantttitle{Mar}{2} \gantttitle{Abr}{2} \gantttitle{May}{2}
  \gantttitle{Jun}{2} \gantttitle{Jul}{2} \\
  \ganttgroup{F1. Análisis}{1}{6} \\
  \ganttbar{Propuesta y estado del arte}{1}{4} \\
  \ganttbar{Requisitos}{4}{6} \\
  \ganttgroup{F2. Diseño e implementación}{6}{15} \\
  \ganttbar{Diseño de la arquitectura}{6}{7} \\
  \ganttbar{Incremento 1: reservas}{8}{10} \\
  \ganttbar{Incremento 2: tiempo real}{11}{13} \\
  \ganttbar{Incremento 3: notificaciones}{14}{15} \\
  \ganttgroup{F3. Validación y cierre}{14}{20} \\
  \ganttbar{Pruebas y validación}{14}{17} \\
  \ganttbar{Redacción de la memoria}{3}{18} \\
  \ganttgroup{Hitos de evaluación}{3}{20} \\
  \ganttmilestone{Informe inicial (sem. 4-6)}{3} \\
  \ganttmilestone{Informe de seguimiento (sem. 10-12)}{6} \\
  \ganttmilestone{Documento final ($\geq$3 sem. antes)}{17} \\
  \ganttmilestone{Depósito ($\geq$5 días hábiles antes)}{19} \\
  \ganttmilestone{Defensa}{20}
\end{ganttchart}

  \caption{Planificación temporal del proyecto (por quincenas).}
  \label{fig:cronograma}
  \fuente{Elaboración propia.}
\end{figure}

\subsection{Desviaciones respecto al plan inicial}
\label{sec:desviaciones}

Compara lo planificado con lo realizado. Por ejemplo: el incremento 2 se
retrasó dos semanas porque la librería de notificaciones elegida no admitía el
modo sin conexión; se sustituyó por otra y se recortó el incremento 3. Si una
tarea planificada no se llegó a hacer (por ejemplo, una prueba con usuarios),
dilo aquí y en el capítulo de pruebas.

\section{Herramientas y tecnologías}
\label{sec:herramientas}

La Tabla~\ref{tab:herramientas} recoge las herramientas utilizadas, con su
versión exacta y el motivo de su elección. El resto de la memoria se refiere a
ellas sin volver a describirlas.

\begin{table}[htbp]
  \centering
  \caption{Herramientas y tecnologías utilizadas.}
  \label{tab:herramientas}
  \small
  \begin{tabularx}{\textwidth}{@{}lllL@{}}
    \toprule
    \textbf{Capa} & \textbf{Herramienta} & \textbf{Versión} & \textbf{Motivo de la elección} \\
    \midrule
    Cliente        & Kotlin~\cite{doc_kotlin}          & 2.0  & Lenguaje recomendado para Android \\
    Servidor       & Python~\cite{doc_python}          & 3.12 & Experiencia previa y ecosistema de pruebas \\
    API            & FastAPI~\cite{doc_fastapi}        & 0.115 & Documentación automática de la \ac{API} \\
    Datos          & PostgreSQL~\cite{doc_postgresql}  & 16   & Transacciones para evitar reservas dobles \\
    Despliegue     & Docker~\cite{doc_docker}          & 27   & Despliegue reproducible con un solo comando \\
    \bottomrule
  \end{tabularx}
  \fuente{Elaboración propia.}
\end{table}

% GUÍA: la versión debe ser la que REALMENTE usaste (mírala en tu proyecto).
% Justifica por adecuación al problema, no por popularidad
% («lo usa la empresa X» no es un motivo técnico).

\section{Viabilidad económica}
\label{sec:viabilidad-economica}

La Tabla~\ref{tab:presupuesto} estima el coste del proyecto. El coste de
personal se calcula con las horas reales de dedicación (un \ac{TFG} de 15 ECTS
supone unas 375 horas) y un coste por hora de referencia; el equipo se imputa
por amortización durante los meses de uso.

\begin{table}[htbp]
  \centering
  \caption{Presupuesto estimado del proyecto (datos de ejemplo).}
  \label{tab:presupuesto}
  \begin{tabular}{@{}lrrr@{}}
    \toprule
    \textbf{Concepto} & \textbf{Cantidad} & \textbf{Coste/ud. (€)} & \textbf{Total (€)} \\
    \midrule
    Horas de desarrollo                    & 375 h    & 20          & \num{7500} \\
    Portátil (amortización 9/48 meses)     & 1        & \num{1200}  & 225 \\
    Móvil de pruebas (amortización 9/36)   & 1        & 300         & 75 \\
    Servidor en la nube                    & 6 meses  & 10          & 60 \\
    Licencias de software                  & ---      & ---         & 0 \\
    \midrule
    \textbf{Total}                         &          &             & \textbf{\num{7860}} \\
    \bottomrule
  \end{tabular}
  \fuente{Elaboración propia.}
\end{table}

\section{Viabilidad legal}
\label{sec:viabilidad-legal}

Indica qué normativa afecta al proyecto y cómo se cumple. Si se tratan datos
personales, aplica el \ac{RGPD}~\cite{rgpd2016} y la ley orgánica española de
protección de datos~\cite{lopdgdd2018}: qué datos se guardan, con qué base
legal, durante cuánto tiempo y cómo se protegen. Revisa también las licencias de
las librerías que usas (y la que eliges para tu código) y, en una aplicación,
las pautas de accesibilidad~\cite{wcag22}.

\section{Sostenibilidad e impacto social}
\label{sec:sostenibilidad}

Valora brevemente el impacto social, económico y ambiental de la solución: a
quién beneficia, qué recursos consume (por ejemplo, energía del servidor o
batería del móvil) y con qué Objetivos de Desarrollo Sostenible se relaciona.
